% Generated by report_gen.py (ai-ee v1) - do not hand-edit.
\documentclass[11pt,a4paper]{article}
\usepackage[margin=2.2cm]{geometry}
\usepackage{graphicx}
\usepackage{booktabs}
\usepackage{longtable}
\usepackage{pdfpages}
\makeatletter
\@ifundefined{KV@Gin@artifact}{\define@key{Gin}{artifact}[]{}}{}
\makeatother
\usepackage{xcolor}
\usepackage[hidelinks]{hyperref}
\setcounter{tocdepth}{1}
\begin{document}
\sloppy
\begin{center}
{\LARGE\bfseries rf-term-150w --- Design Document}\\[6pt]
{\large ai-ee v1 pipeline}\\[2pt]
generated 2026-08-09 04:13:59
\end{center}
\tableofcontents

\section{Board and Run Metadata}
{\small
\begin{longtable}{lp{11cm}}
\toprule
\textbf{Field} & \textbf{Value} \\
\midrule
\endhead
board & rf-term-150w \\
workspace & \texttt{boards/rf-term-150w} \\
phase & P6 \\
gate status & all 6 recorded gates pass \\
state created & 2026-08-08T22:33:40 \\
state updated & 2026-08-09T04:12:24 \\
\bottomrule
\end{longtable}
}
\section{Overview}
\subsection*{Standing assumptions (user said: do not ask, assume and record)}
These are orchestrator-level assumptions handed to every downstream agent. Every one of them must appear in the final README.

A1. \textbf{Band}: DC - 25 MHz. Fc = 25 MHz taken as the worst-case design point (highest residual-reactance mismatch). "Fc +/-10\%" = 22.5 - 27.5 MHz; the design targets \textgreater{}= 20 dB from DC to 27.5 MHz, a superset.

A2. \textbf{Duty (left blank by the user)}: assume \textbf{CW, 100\% duty, continuous}. This is the thermally worst case; a pulsed spec can only be easier. All thermal numbers in the README are steady-state CW at 150 W.

A3. \textbf{Ambient / cooling}: 25 C ambient, natural convection for the "derated air-cooled power" figure; forced-air / coldplate handled by quoting the required heatsink thermal resistance instead.

A4. \textbf{Mechanical}: the resistor bolts DIRECTLY to the user's heatsink or coldplate through its own flange holes. The PCB is a launch/interface carrier only, sits on the same surface, and is not in the thermal path. No heat is intentionally conducted through FR4.

A5. \textbf{Assembly}: hand-solder / bench build of 5. Not JLC PCBA (the flange resistor and the SMA jack are through-hole/mechanical parts; PCBA of a 5-off with a bolt-down power part is not sensible). BOM+CPL are still produced as deliverables.

A6. \textbf{Adjustment element}: a shunt trimmer capacitor at the launch is the operator-adjustable reactance null. It must be reachable by a tuning tool with the board bolted down and the SMA cable mated (adjust from the top, no desoldering, no unbolting).

A7. \textbf{RF power vs DC test}: "DC resistance 50 ohm +/-2\%" is met by the resistor element's own tolerance (+/-1\% or +/-2\% part). No trimming of the DC value is provided or required.

A8. \textbf{Cost}: the \$40 / build-of-5 cap is BOM + bare PCB. PCB shipping and any tax are excluded (they are order-time, not design, costs) - if the cap still cannot be met, report the real number rather than degrading the design.

A9. \textbf{No ordering}: deliverables stop at fab artifacts (gerbers, drill, BOM, CPL) + README. The pipeline runs P0-P9; P10/ordering is not requested. \textless{}/content\textgreater{} \textless{}/invoke\textgreater{}

\subsection*{Stackup}
Chosen stackup: JLC2313\_1.6 (2 layer, 1.6 mm, 1 oz outer, HASL)
\section{Requirements}
\subsection*{requirements.md - rf-term-150w}
Single-port 50 ohm / 150 W CW RF termination (dummy load head), DC - 25 MHz.

Source: \texttt{brief/user-brief.md} (user's words, verbatim) + \texttt{brief/brief.md} (orchestrator standing assumptions A1-A9). The user instructed: \textbf{do not ask clarifying questions - assume and record.} Every gap is therefore closed inline as \texttt{ASSUMED:} and section 9 is \texttt{none}.

\textbf{A1-A9 are decided answers, not proposals.} They are reproduced in section 11 for traceability and must all appear in the final README.

---

\subsubsection*{1. Function}
A passive, single-port 50 ohm RF termination head. One SMA female jack accepts up to 150 W of continuous-wave RF power from DC to 25 MHz over a standard male-terminated coax cable and dissipates all of it as heat in a bolt-down termination element. The board itself is the RF launch and the mechanical interface only: it carries the connector, the resistive element's electrical connection, and one operator-adjustable reactance null. It is NOT in the thermal path - the termination element bolts through its own flange directly to an external heatsink or coldplate that the user supplies and that is out of scope for this design (A4). There is no active circuitry, no DC supply, no control, no telemetry, and no signal path other than the single RF port.

Design point (A1): Fc = 25 MHz taken as worst case (highest residual-reactance mismatch); the DC-25 MHz band means every frequency below 25 MHz is easier. Duty is CW / 100\% (A2).

\paragraph*{Derived electrical operating point (arithmetic, stated per assignment)}
At 150 W into 50 ohm:

\begin{quote}\small\ttfamily\raggedright
\textbar{} Quantity \textbar{} Formula \textbar{} Value \textbar{}\\
\textbar{}---\textbar{}---\textbar{}---\textbar{}\\
\textbar{} Port voltage, rms \textbar{} V = sqrt(P{*}R) = sqrt(150 {*} 50) = sqrt(7500) \textbar{} {*}{*}86.6 Vrms{*}{*} \textbar{}\\
\textbar{} Port voltage, peak \textbar{} V {*} sqrt(2) \textbar{} {*}{*}122.5 Vpeak{*}{*} \textbar{}\\
\textbar{} Port voltage, pk-pk \textbar{} 2 {*} Vpeak \textbar{} {*}{*}244.9 Vpp{*}{*} \textbar{}\\
\textbar{} Port current, rms \textbar{} I = sqrt(P/R) = sqrt(3) \textbar{} {*}{*}1.732 Arms{*}{*} \textbar{}\\
\textbar{} Port current, peak \textbar{} I {*} sqrt(2) \textbar{} {*}{*}2.449 Apeak{*}{*} \textbar{}
\end{quote}
Consequences that bind every later phase:

\begin{itemize}
\item \textbf{Any component at the port node sees 86.6 Vrms / 122.5 Vpeak / 244.9 Vpp continuously.}
That includes the adjustment element of A6. A 100 V-rated trimmer capacitor is DISQUALIFIED.
Minimum acceptable rating: \textbf{\textgreater{}= 250 V working}, and the part must be rated for continuous RF
(not just DC) at 25 MHz. See section 8 flag F1.
\item Current is 1.732 Arms - \textbf{below the 3 A high-current threshold}, so the high-current
compliance flag does NOT apply (section 8). Conductor sizing is driven by the voltage/clearance
case and by skin effect at 25 MHz (skin depth in Cu at 25 MHz = 13.1 um; 1 oz = 35 um is
\textasciitilde{}2.7 skin depths), not by DC ampacity.
\item \textbf{Electrically small board.} Free-space wavelength at 25 MHz is 12 m; in FR4 microstrip
(eps\_eff \textasciitilde{} 3.3) it is \textasciitilde{}6.6 m. The 30 mm outline is \textasciitilde{}0.005 lambda. Transmission-line behaviour is
negligible and the design is purely lumped. This is why the brief's "no controlled-impedance
service" constraint costs nothing - it is not a risk. Record it; do not spend budget on it.
\end{itemize}
\paragraph*{Derived match budget (arithmetic)}
\begin{quote}\small\ttfamily\raggedright
\textbar{} Spec \textbar{} Reflection coefficient \textbar{} VSWR \textbar{} Max residual series reactance \textbar{}\\
\textbar{}---\textbar{}---\textbar{}---\textbar{}---\textbar{}\\
\textbar{} RL \textgreater{}= 26 dB at 25 MHz \textbar{} \textbackslash{}\textbar{}G\textbackslash{}\textbar{} \textless{}= 10\textasciicircum{}(-26/20) = {*}{*}0.0501{*}{*} \textbar{} \textless{}= {*}{*}1.106{*}{*} \textbar{} \textbackslash{}\textbar{}X\textbackslash{}\textbar{} \textless{}= {*}{*}5.0 ohm{*}{*} = 31.9 nH at 25 MHz \textbar{}\\
\textbar{} RL \textgreater{}= 20 dB, 22.5-27.5 MHz \textbar{} \textbackslash{}\textbar{}G\textbackslash{}\textbar{} \textless{}= {*}{*}0.100{*}{*} \textbar{} \textless{}= {*}{*}1.222{*}{*} \textbar{} \textbackslash{}\textbar{}X\textbackslash{}\textbar{} \textless{}= {*}{*}10.05 ohm{*}{*} = 71.1 nH at 22.5 MHz / 58.2 nH at 27.5 MHz \textbar{}
\end{quote}
Derivation: for Z = R + jX with R = 50, \textbackslash{}\textbar{}G\textbackslash{}\textbar{} = X / sqrt(100\textasciicircum{}2 + X\textasciicircum{}2), so \textbackslash{}\textbar{}G\textbackslash{}\textbar{} = 0.0501 -\textgreater{} X = 5.02 ohm. The +/-2\% DC resistance tolerance alone costs \textbackslash{}\textbar{}G\textbackslash{}\textbar{} = 1/101 = 0.0099 (RL 40 dB); taken in quadrature it leaves 4.92 ohm of reactive budget instead of 5.02 - i.e. \textbf{\textasciitilde{}0.1 dB}. Resistance tolerance is NOT the binding constraint; \textbf{residual series inductance is}. The whole point of the adjustment element is to cancel it.

\textbf{Post-tune residual series inductance must land under \textasciitilde{}32 nH at 25 MHz.} For scale: \textasciitilde{}20 nH per 25 mm of thin wire/lead. The mechanical layout of the launch is therefore the primary performance lever, and the trimmer only has to absorb the leftovers.

\subsubsection*{2. Interfaces}
\begin{quote}\small\ttfamily\raggedright
\textbar{} \# \textbar{} Interface \textbar{} Direction \textbar{} Standard / detail \textbar{} Connector \textbar{}\\
\textbar{}---\textbar{}---\textbar{}---\textbar{}---\textbar{}---\textbar{}\\
\textbar{} I1 \textbar{} RF port \textbar{} in (absorbs) \textbar{} 50 ohm nominal, DC - 25 MHz, up to 150 W CW. 86.6 Vrms / 1.732 Arms. DC-continuous to ground through the termination element (the "DC resistance 50 ohm +/-2\%" spec is measured here). \textbar{} {*}{*}SMA female (jack){*}{*}, stated in brief; mates a standard SMA male cable end \textbar{}\\
\textbar{} I2 \textbar{} Tuning access \textbar{} mechanical \textbar{} Operator-adjustable reactance null, reachable by a hand tuning tool {*}{*}from the top, with the board bolted down and the SMA cable mated{*}{*} - no desoldering, no unbolting (A6, brief). \textbar{} trimmer screw / slot \textbar{}\\
\textbar{} I3 \textbar{} Thermal / mechanical ground \textbar{} mechanical \textbar{} Termination element flange bolts DIRECTLY to the user's heatsink or coldplate through its own flange holes (A4). PCB sits on the same surface, not in the thermal path. \textbar{} user-supplied fasteners, out of scope \textbar{}
\end{quote}
Notes binding the interfaces:

\begin{itemize}
\item \textbf{This is the only electrical interface on the board.} No DC input, no bias tee, no monitor/
detector tap, no LEDs, no test points are called for. \texttt{ASSUMED:} none wanted - the brief's
\textless{}= 4 BOM lines / \textless{}= 6 placements budget forbids scope creep anyway.
\item \texttt{ASSUMED:} \textbf{RF ground reference is the SMA jack body plus the ground pour}, tied to the
termination element's cold end. No separate ground lug or chassis screw is specified. Bonding
to the heatsink happens incidentally through the resistor flange and any mounting hardware.
\item \texttt{ASSUMED:} \textbf{SMA jack is a through-hole, PCB-mount part} (edge-launch or vertical), hand-
solderable per A5. Exact style (edge vs vertical vs 4-hole flange) is left to P1/P2 parts
selection; the binding constraint is that the ground return around the launch be short and
wide, since launch inductance eats the 32 nH budget of section 1.
\item \textbf{SMA power handling check (record, do not re-open):} SMA is rated \textasciitilde{}335 Vrms breakdown; at
86.6 Vrms it has \textasciitilde{}3.9x voltage margin, and at 25 MHz there is no meaningful connector loss.
SMA is electrically adequate for 150 W CW here. The real SMA caveat is mechanical - mating-cycle
wear and centre-pin heating under repeated hot connects - which is a usage note for the README,
not a design change (the user specified SMA).
\item \textbf{No isolation, no protection.} An open or shorted port, or drive above 150 W, is outside the
design envelope; nothing on the board detects or survives it. README must say so.
\end{itemize}
\subsubsection*{3. Power}
\textbf{There is no DC supply and there are no power rails on this board.} It is fully passive.

The only "power" is the RF drive absorbed at I1:

\begin{quote}\small\ttfamily\raggedright
\textbar{} Param \textbar{} Value \textbar{} Status \textbar{}\\
\textbar{}---\textbar{}---\textbar{}---\textbar{}\\
\textbar{} RF input power \textbar{} 150 W {*}{*}CW, 100\% duty, continuous{*}{*} \textbar{} A2 (user left duty blank; CW is the thermal worst case - any pulsed spec is strictly easier) \textbar{}\\
\textbar{} Frequency range \textbar{} DC - 25 MHz; design point 25 MHz \textbar{} A1 \textbar{}\\
\textbar{} Port voltage / current \textbar{} 86.6 Vrms / 1.732 Arms (section 1) \textbar{} derived \textbar{}\\
\textbar{} Power dissipated on-board \textbar{} \textasciitilde{}150 W, essentially all of it in the termination element \textbar{} derived \textbar{}\\
\textbar{} Power dissipated in the PCB \textbar{} intended \textasciitilde{}0 W. No heat is intentionally conducted through FR4 (A4) \textbar{} A4 \textbar{}\\
\textbar{} Battery / charging \textbar{} {*}{*}none{*}{*} \textbar{} N/A \textbar{}\\
\textbar{} Mains \textbar{} {*}{*}none{*}{*} \textbar{} N/A \textbar{}
\end{quote}
Secondary dissipation to budget at P2/P3 (small but not zero):

\begin{itemize}
\item The adjustment element carries RF current at the port node. For a shunt trimmer of value C,
I = V {*} 2\emph{pi}f{*}C; e.g. 10 pF at 25 MHz and 86.6 Vrms -\textgreater{} 0.136 Arms, and with Q = 200
(Xc = 637 ohm, ESR = 3.2 ohm) that is \textasciitilde{}60 mW. Small - but it means the trimmer's \textbf{voltage}
rating, not its power rating, is the gate. Re-check against the actual chosen part at P2.
\item Connector and solder-joint I\textasciicircum{}2R at 1.73 Arms is negligible (\textless{}\textless{} 1 W) but the joints must not be
the mechanical support for a heavy flange part.
\end{itemize}
\paragraph*{Thermal requirement derived from the derating curve (drives part selection at P1)}
The brief's README deliverable demands the required heatsink thermal resistance computed {*}{*}from the actual datasheet derating curve{*}{*}. The arithmetic and its consequence are a requirement, not an implementation detail, because it determines whether a candidate part is admissible at all:

T\_flange\_allowed = T\_max - (P\_op / P\_rated) {*} (T\_max - T\_ref) Rth\_sink\_max     = (T\_flange\_allowed - T\_ambient) / P\_op        {[}includes the interface{]}

With the \textbf{generic} flange-resistor derating curve (\texttt{ASSUMED:} P\_rated at T\_ref = 25 C flange, derating linearly to zero at T\_max = 150 C), T\_ambient = 25 C (A3), P\_op = 150 W:

\begin{quote}\small\ttfamily\raggedright
\textbar{} Part's rated power \textbar{} Allowed flange temp at 150 W \textbar{} Max total Rth (sink + interface) \textbar{}\\
\textbar{}---\textbar{}---\textbar{}---\textbar{}\\
\textbar{} 150 W \textbar{} 25 C \textbar{} {*}{*}0.00 C/W - IMPOSSIBLE{*}{*} \textbar{}\\
\textbar{} 250 W \textbar{} 75 C \textbar{} 0.333 C/W \textbar{}\\
\textbar{} 300 W \textbar{} 87.5 C \textbar{} 0.417 C/W \textbar{}\\
\textbar{} 400 W \textbar{} 103.1 C \textbar{} 0.521 C/W \textbar{}
\end{quote}
{*}{*}Binding conclusion: a part rated exactly 150 W cannot dissipate 150 W with ANY finite heatsink at 25 C ambient.\textbf{ The termination element must be rated }\textgreater{}= 250 W{*}{*} at its datasheet reference flange temperature (or carry an equivalent curve that leaves a non-zero allowed flange rise at 150 W). P1 MUST re-derive this table from the real datasheet curve of the part actually chosen and put the result in the README; the table above is a screening filter, not the answer.

Consequence for A3's "derated air-cooled power" figure: the same equation solved for natural convection at 25 C ambient with a realistic bare-flange Rth. Expect it to be a small fraction of 150 W - i.e. \textbf{the board is unusable at rated power without the external heatsink}, which the README must state plainly.

\subsubsection*{4. Environment}
\begin{quote}\small\ttfamily\raggedright
\textbar{} Item \textbar{} Value \textbar{} Status \textbar{}\\
\textbar{}---\textbar{}---\textbar{}---\textbar{}\\
\textbar{} Ambient temperature \textbar{} {*}{*}25 C{*}{*} for all thermal numbers \textbar{} A3 \textbar{}\\
\textbar{} Cooling \textbar{} natural convection for the "derated air-cooled power" figure; forced air / coldplate covered by quoting the required Rth instead \textbar{} A3 \textbar{}\\
\textbar{} Enclosure \textbar{} {*}{*}none{*}{*} - open bench / bolted to a user heatsink \textbar{} `ASSUMED:` (brief describes a bolt-on heatsink or coldplate, no enclosure) \textbar{}\\
\textbar{} Ingress protection \textbar{} {*}{*}none required{*}{*} \textbar{} `ASSUMED:` indoor bench use \textbar{}\\
\textbar{} Vibration / shock \textbar{} {*}{*}none{*}{*} - stationary bench use \textbar{} `ASSUMED:` \textbar{}\\
\textbar{} Altitude / humidity \textbar{} ordinary indoor lab conditions \textbar{} `ASSUMED:` \textbar{}\\
\textbar{} Hot surfaces \textbar{} the termination element and its flange WILL run far above touch-safe temperature at 150 W CW - see section 8 flag F3 \textbar{} derived \textbar{}
\end{quote}
\texttt{ASSUMED:} no storage/operating temperature range beyond the above was stated; standard FR4 (Tg 130-150 C) and standard-tolerance passives are acceptable. Note the PCB sits on the same surface as a flange that may exceed 100 C, so the board WILL be heated by proximity even though it is not intentionally in the thermal path (A4). P6 should keep FR4 and any polymer-bodied part away from the flange footprint where the outline allows.

\subsubsection*{5. Size \& mounting}
\begin{quote}\small\ttfamily\raggedright
\textbar{} Item \textbar{} Value \textbar{} HARD / soft \textbar{}\\
\textbar{}---\textbar{}---\textbar{}---\textbar{}\\
\textbar{} Board outline \textbar{} {*}{*}\textless{}= 30 x 30 mm{*}{*} \textbar{} {*}{*}HARD{*}{*} - stated as a constraint in the brief. Binds permanently at P5 `board\_init`; cannot be relaxed later without restarting layout. \textbar{}\\
\textbar{} Layer count \textbar{} {*}{*}2 layers{*}{*}, JLCPCB standard FR4 process only. No upcharge options, no controlled-impedance service. \textbar{} {*}{*}HARD{*}{*} - stated in the brief \textbar{}\\
\textbar{} PCB thickness \textbar{} 1.6 mm \textbar{} `ASSUMED:` JLCPCB standard, no upcharge (brief forbids upcharge options) \textbar{}\\
\textbar{} Copper weight \textbar{} 1 oz outer \textbar{} `ASSUMED:` JLCPCB standard; adequate - see the skin-depth note in section 1 \textbar{}\\
\textbar{} Mounting holes \textbar{} none specified \textbar{} `ASSUMED:` the termination element's own flange screws are the primary mechanical anchor (A4). Add PCB clearance/mounting holes only as the mechanical stack-up demands. \textbar{}\\
\textbar{} Height limit \textbar{} none stated \textbar{} `ASSUMED:` unconstrained. The tallest parts will be the SMA jack and the trimmer, both of which must stay accessible per I2. \textbar{}
\end{quote}
Mechanical items that later phases MUST resolve (recorded here because they can invalidate the HARD outline):

1. \textbf{Flange bolt pattern vs the 30 x 30 mm outline.} A \textgreater{}= 250 W flange resistor is a physically large part; its flange holes bolt to the heatsink, not to the PCB. The outline must therefore accommodate a notch, cutout, or edge-relief for the flange and its screws, alongside an SMA jack, inside 30 x 30 mm. \texttt{ASSUMED:} this fits; {*}{*}if P1's chosen part makes it impossible, that is a HARD-constraint conflict and must be escalated, not silently resized.{*}{*} 2. \textbf{Terminal-height stack-up.} The PCB sits on the same surface as the flange (A4), so the board's top copper sits \textasciitilde{}1.6 mm above the flange plane while the resistor's terminals sit at or near flange level. \texttt{ASSUMED:} the terminal tabs are formed/bridged to the PCB pads with a short solder connection or strap. Whatever bridges that gap adds series inductance against the 32 nH budget of section 1 - \textbf{keep it short and wide.} Resolve concretely at P2/P6. 3. \textbf{Tuning-tool access cone.} Per I2/A6 the adjustment must be reachable from the top with the SMA cable mated and the board bolted down. That forbids placing the trimmer under the SMA barrel or hard against the flange body. This is a placement constraint at P6, not a preference.

\subsubsection*{6. Quantity \& budget}
\begin{quote}\small\ttfamily\raggedright
\textbar{} Item \textbar{} Value \textbar{} Status \textbar{}\\
\textbar{}---\textbar{}---\textbar{}---\textbar{}\\
\textbar{} Build quantity \textbar{} {*}{*}5{*}{*} \textbar{} brief \textbar{}\\
\textbar{} Total budget \textbar{} {*}{*}\textless{}= \$40 for the build of 5{*}{*} = {*}{*}\$8.00 per board{*}{*}, covering BOM at qty-5 pricing + bare PCB \textbar{} brief \textbar{}\\
\textbar{} Excluded from the cap \textbar{} PCB shipping and tax - order-time costs, not design costs \textbar{} A8 \textbar{}\\
\textbar{} Sourcing \textbar{} all parts {*}{*}in stock at LCSC or DigiKey{*}{*}; no NRND, no obsolete \textbar{} brief (HARD) \textbar{}\\
\textbar{} Cap behaviour \textbar{} if the cap cannot be met, {*}{*}report the real number - do not degrade the design{*}{*} \textbar{} A8 \textbar{}
\end{quote}
\textbf{KNOWN BUDGET RISK - record now, resolve at P1, do not paper over.} Section 3 concludes the termination element must be rated \textgreater{}= 250 W. Flange-mount RF terminations in that class routinely cost \$20-60 each in single quantities, i.e. {*}{*}one part could exceed the entire \$40 build-of-5 cap by several times.{*}{*} Levers available to P1, in preference order:

1. A non-RF-specific bolt-down thick-film power resistor (TO-220/TO-247/flanged) of the right value and power class. Admissible here \emph{only because} the +/-10 ohm reactance budget of section 1 plus the adjustment element of A6 tolerate a part that is not characterised for RF - at 25 MHz on an electrically tiny board that is a real option, whereas at 500 MHz it would not be. P1 must verify parasitic inductance against the 32 nH budget before committing. 2. Accept the overrun and report the true number per A8.

Do \textbf{not} resolve this by dropping the power rating back to 150 W - section 3 proves that is thermally unsatisfiable.

\subsubsection*{7. Assembly}
\begin{quote}\small\ttfamily\raggedright
\textbar{} Item \textbar{} Value \textbar{} Status \textbar{}\\
\textbar{}---\textbar{}---\textbar{}---\textbar{}\\
\textbar{} Method \textbar{} {*}{*}hand-solder / bench build of 5.{*}{*} NOT JLC PCBA - the flange resistor and the SMA jack are through-hole/mechanical parts, and PCBA of a 5-off with a bolt-down power part is not sensible \textbar{} A5 \textbar{}\\
\textbar{} Deliverables regardless \textbar{} BOM and CPL are still produced as fab deliverables \textbar{} A5, brief \textbar{}\\
\textbar{} Assembly sides \textbar{} {*}{*}single-sided (top) only{*}{*} \textbar{} `ASSUMED:` derived from A4 - the bottom face sits flat on the heatsink surface, so it must carry no components \textbar{}\\
\textbar{} Unique BOM lines \textbar{} {*}{*}\textless{}= 4{*}{*} \textbar{} brief (HARD) \textbar{}\\
\textbar{} Total placements \textbar{} {*}{*}\textless{}= 6{*}{*} \textbar{} brief (HARD) \textbar{}\\
\textbar{} Layers \textbar{} 2, JLCPCB standard FR4, no upcharge options \textbar{} brief (HARD) \textbar{}
\end{quote}
Budget headroom note: the function needs roughly a jack + a termination element + an adjustment element = \textasciitilde{}3 lines / \textasciitilde{}3 placements. The caps are not tight; they exist to forbid scope creep. Mounting holes and cutouts are not placements and do not consume BOM lines.

\subsubsection*{8. Compliance / safety flags}
Flags that \textbf{APPLY}:

\textbf{F1. \textgreater{}30 V present - 86.6 Vrms / 122.5 Vpeak / 244.9 Vpp at the RF port, continuously.} This is well above the 30 V flag threshold and above the 60 VDC / 30 Vrms conventional touch-safe limit. Consequences that bind later phases:

\begin{itemize}
\item Every component on the port node - explicitly including the A6 adjustment trimmer - must be
rated for continuous service at \textgreater{}= 250 V working. \textbf{A 100 V-rated trimmer is disqualified.}
\item Creepage/clearance on the port net must be sized for 245 Vpp, not for a signal-level net.
This is a 2-layer board with an exposed top side; a trimmer screw the operator touches while
the port is live is a live-part-access question, so the tuning access of I2 must be designed to
keep fingers off the hot node (insulated tuning tool, grounded rotor orientation).
\item \textbf{Tuning while transmitting at 150 W is the intended use case} (adjust for minimum reflected
power). The README's tuning procedure must therefore state the shock precaution explicitly, or
prescribe tuning at reduced drive power.
\end{itemize}
\textbf{F2. RF transmit / high-power RF at the port - 150 W CW, DC-25 MHz.} The board is a \emph{load}, not a source, so it does not itself radiate by design; but it is permanently connected to a 150 W HF transmitter and it is unshielded and unenclosed.

\begin{itemize}
\item A well-matched termination radiates little, but any residual mismatch, an unmated port, or a
poorly made cable connection turns the assembly into an accidental HF radiator at 150 W. The
README must state: \textbf{never key the transmitter with the port unmated.}
\item 25 MHz sits just below the 27.12 MHz ISM allocation and inside HF spectrum used by licensed
services. This assembly is a test load, not a certifiable product; it is not intended for
connection to an antenna.
\item RF burn / RF exposure hazard exists at the port at 150 W independently of the thermal hazard
below. Do not handle the centre conductor or an unmated connector while keyed.
\end{itemize}
\textbf{F3. Burn / touch-temperature hazard - the intended, normal operating condition.} Section 3 shows the flange must sit at 75-103 C for the derating curve to close at 150 W, and the resistor's own body will be hotter still. The heatsink it bolts to will also be hot. This is not a fault condition; it is what "working correctly" looks like.

\begin{itemize}
\item Requires a "HOT SURFACE" warning in the README and, budget permitting within the 30 x 30 mm
outline, on the silkscreen.
\item Cool-down time before handling must be stated in the README.
\item The PCB (FR4, and any polymer-bodied part on it) sits on the same hot surface - see section 4.
\end{itemize}
Flags that do \textbf{NOT} apply, stated explicitly:

\begin{itemize}
\item \textbf{Mains voltage: none.} Nothing on this board connects to mains. The user's brief fixed this.
\item \textbf{Battery: none.} No cell, no charger, no chemistry decision anywhere.
\item \textbf{Motors: none.}
\item \textbf{High current (\textgreater{}3 A): does not apply.} Port current is 1.732 Arms (section 1), below the
threshold.
\end{itemize}
\textbf{Why these are documented and not asked:} the user's brief already fixed the safety-relevant architecture - external heatsink supplied by the user, no mains, no battery, passive board. The normal P0 rule "ask, never guess, for anything in section 8" is satisfied because nothing in section 8 is \emph{guessed}: F1 and F2 follow arithmetically from the user's own 150 W / 50 ohm / 25 MHz numbers, and F3 follows from the user's own external-heatsink architecture. There is no open safety decision left to make, so there is nothing safety-blocking to ask.

\subsubsection*{9. Open questions}
\textbf{none.}

The user instructed: "Do not ask clarifying questions. Make assumptions, record them in the README, and proceed." Every gap is closed as an \texttt{ASSUMED:} above or by standing assumption A1-A9 (section 11). Nothing remaining is safety-blocking - see the closing note in section 8.

The three assumptions with the largest downstream leverage, flagged for the README rather than asked: (a) CW duty per A2 - if the real use is pulsed, the whole thermal case relaxes; (b) the generic derating curve in section 3 - P1 must replace it with the real datasheet curve; (c) the mechanical stack-up items in section 5, which are the only route by which the HARD 30 x 30 mm outline could be found infeasible.

\subsubsection*{10. Acceptance criteria (measurable, testable)}
Pass/fail list for the finished board. Each is measurable with bench equipment or from the fab artifacts.

1. \textbf{DC resistance.} Measured port centre-to-shell DC resistance = \textbf{50 ohm +/-2\%}, i.e. \textbf{49.00 to 51.00 ohm} inclusive. Measured with a 4-wire ohmmeter at the mated SMA interface, at room temperature, unpowered. Met by the termination element's own tolerance; {*}{*}no DC trimming is provided or required{*}{*} (A7). 2. \textbf{Return loss at the design point.} \textbf{\textgreater{}= 26 dB at 25.0 MHz} after adjustment (\textbackslash{}\textbar{}G\textbackslash{}\textbar{} \textless{}= 0.0501, VSWR \textless{}= 1.106). Measured on a calibrated VNA at the SMA port with the board bolted to its heatsink and the trimmer set per the documented procedure. 3. \textbf{Return loss across the band.} \textbf{\textgreater{}= 20 dB at every frequency from 22.5 to 27.5 MHz} (\textbackslash{}\textbar{}G\textbackslash{}\textbar{} \textless{}= 0.100, VSWR \textless{}= 1.222), with the trimmer left at its criterion-2 setting - i.e. one adjustment satisfies both 2 and 3. Per A1 the design targets this from DC to 27.5 MHz, a superset of the specified window; verify the superset, report the specified window. 4. \textbf{Residual reactance (diagnostic for 2 and 3).} Post-tune equivalent series reactance at 25 MHz \textbf{\textbackslash{}\textbar{}X\textbackslash{}\textbar{} \textless{}= 5.0 ohm} (equivalently \textless{}= 31.9 nH residual series inductance), and \textbf{\textbackslash{}\textbar{}X\textbackslash{}\textbar{} \textless{}= 10.05 ohm} at both band edges. This is the internal budget that criteria 2 and 3 test from outside; record it so a failure is diagnosable. 5. \textbf{Adjustment range, documented and demonstrated.} The trimmer must move the port reactance at 25 MHz over a \textbf{continuous range that spans at least 0 to 10 ohm} (0 to \textasciitilde{}64 nH equivalent) so it can absorb the as-built residual with margin at both ends. The achieved range must be \textbf{measured and stated as a number in the README}, not asserted. 6. \textbf{Adjustment is operator-accessible as specified.} Demonstrated: the trimmer can be turned through its full range with a hand tuning tool while the board is bolted down AND the SMA cable is mated, \textbf{without desoldering and without unbolting anything} (A6, I2). 7. \textbf{Adjustment-element voltage rating.} Every component on the port node is rated \textbf{\textgreater{}= 250 V working}, verified against its datasheet, versus the computed 122.5 Vpeak / 244.9 Vpp. Checked from the BOM, not the bench. 8. \textbf{BOM line count.} \textbf{\textless{}= 4 unique BOM line items.} Counted from the delivered BOM file. 9. \textbf{Placement count.} \textbf{\textless{}= 6 total placements.} Counted from the delivered CPL file. 10. \textbf{Outline.} Board fits within \textbf{30.0 x 30.0 mm} - HARD. Verified from the delivered gerbers / board edge-cut extents. 11. \textbf{Process.} \textbf{2 layers, JLCPCB standard FR4}, no upcharge options, no controlled-impedance service on the order. Verified from the fab notes and the delivered stackup. 12. \textbf{Cost.} \textbf{\textless{}= \$40.00 total for a build of 5}, computed as BOM at qty-5 pricing + bare PCB, excluding shipping and tax (A8). If exceeded, the \textbf{actual} figure is reported and the design is NOT degraded to meet the cap. 13. \textbf{Sourcing.} Every BOM line is \textbf{in stock at LCSC or DigiKey} at the time of the BOM, and no line is NRND or obsolete. Verified per line with a stock figure and a date. 14. \textbf{Thermal documentation (README deliverable).} The README states, computed from the {*}{*}actual datasheet derating curve{*}{*} of the part actually chosen: (a) the allowed flange temperature at 150 W, (b) the required total heatsink-plus-interface thermal resistance to hold it at 25 C ambient, and (c) the \textbf{derated air-cooled power} achievable with no heatsink under natural convection at 25 C. All three as numbers with the curve they came from. 15. \textbf{Gate cleanliness.} ERC clean and DRC clean against the JLCPCB 2-layer rule set. Verified by the pipeline gates, not by inspection. 16. \textbf{Fab artifacts complete.} Gerbers, drill, BOM, CPL, plus the one-page README containing the tuning procedure, the measured tuning range (criterion 5), and the thermal numbers (criterion 14). Per A9, the pipeline stops at fab artifacts; \textbf{no order is placed.}

\subsubsection*{11. Traceability - standing assumptions A1-A9}
Reproduced verbatim in substance from \texttt{brief/brief.md}. These are \textbf{decided answers}. All of them must appear in the final README.

\begin{quote}\small\ttfamily\raggedright
\textbar{} \# \textbar{} Assumption \textbar{} Where it binds \textbar{}\\
\textbar{}---\textbar{}---\textbar{}---\textbar{}\\
\textbar{} A1 \textbar{} Band DC - 25 MHz; Fc = 25 MHz as worst case; Fc +/-10\% = 22.5-27.5 MHz; design targets \textgreater{}= 20 dB DC-27.5 MHz (a superset) \textbar{} sections 1, 10 (criteria 2, 3) \textbar{}\\
\textbar{} A2 \textbar{} Duty (left blank by user) = {*}{*}CW, 100\%, continuous{*}{*} - thermal worst case; a pulsed spec can only be easier \textbar{} sections 3, 4 \textbar{}\\
\textbar{} A3 \textbar{} 25 C ambient; natural convection for the "derated air-cooled power" figure; forced air / coldplate handled by quoting required Rth \textbar{} sections 3, 4, 10 (criterion 14) \textbar{}\\
\textbar{} A4 \textbar{} Resistor bolts DIRECTLY to the user's heatsink/coldplate through its own flange holes; PCB is a launch/interface carrier on the same surface, not in the thermal path; no heat intentionally through FR4 \textbar{} sections 1, 3, 4, 5, 7 \textbar{}\\
\textbar{} A5 \textbar{} Hand-solder / bench build of 5, not JLC PCBA; BOM + CPL still produced \textbar{} section 7 \textbar{}\\
\textbar{} A6 \textbar{} Adjustment element = shunt trimmer capacitor at the launch; reachable by a tuning tool with the board bolted down and the SMA cable mated (top access, no desoldering, no unbolting) \textbar{} sections 2, 5, 8, 10 (criteria 5, 6) \textbar{}\\
\textbar{} A7 \textbar{} "DC resistance 50 ohm +/-2\%" met by the element's own tolerance (+/-1\% or +/-2\% part); no DC trimming provided or required \textbar{} section 10 (criterion 1) \textbar{}\\
\textbar{} A8 \textbar{} \$40 / build-of-5 cap = BOM + bare PCB; shipping and tax excluded; if the cap cannot be met, report the real number rather than degrading the design \textbar{} sections 6, 10 (criterion 12) \textbar{}\\
\textbar{} A9 \textbar{} No ordering. Deliverables stop at fab artifacts (gerbers, drill, BOM, CPL) + README. Pipeline runs P0-P9; P10/ordering not requested \textbar{} section 10 (criterion 16) \textbar{}
\end{quote}
\section{Architecture}
\subsection*{Decisions of record (state.json)}
{\small
\begin{longtable}{lp{5.2cm}p{7.2cm}}
\toprule
\textbf{Phase} & \textbf{What} & \textbf{Why} \\
\midrule
\endhead
P0 & Fc = 25 MHz worst case; band DC-25 MHz; spec band 22.5-27.5 MHz & user answered 'up to 25 MHz' mid-run; top of band is the worst case for residual-reactance mismatch \\
P0 & Duty assumed CW 100 percent continuous at 150 W & user left the duty field blank and forbade clarifying questions; CW is the thermally bounding case \\
P0 & Spawn tier substitution: fable roles run on opus for this build & fable 5 returned out-of-usage-credits; playbook allows escalating a tier, never silent downgrade \\
P1 & Board-wide copper clearance set to 0.80 mm & 150W into 50R = 86.6 Vrms / 122.5 Vpeak -\textgreater{} IPC-2221 101-150 V band. Exposed HASL lands take row A6 = 0.80 mm (masked traces would be B4 = 0.40 mm). Board is nearly empty so the strictest row is free; avoids per-item creepage adjudication entirely (LEARNINGS 2026-07-29 check\_creepage). \\
P1 & Stackup JLC2313\_1.6 (2L, 1.6mm, 1oz HASL); no controlled impedance & only non-upcharge 2-layer stackup modeled; at 25 MHz a 15 mm launch is 0.8 degrees electrical (0.002 lambda) so trace Z0 is irrelevant - the no-controlled-impedance constraint costs nothing \\
P1 & Flange resistor gets an EDGE NOTCH, not an interior window & board\_init rejects interior cutouts by design - an inner Edge.Cuts rect silently becomes the board outline (LEARNINGS 2026-07-28 geom/layout). Notch also lets the resistor bolt straight to the heatsink. \\
P0 & Termination element must be rated \textgreater{}= 250 W, not 150 W & P0 arithmetic: a part rated 150 W at 25 C flange derating to zero at 150 C needs Rth\_heatsink = 0 C/W to pass 150 W at 25 C ambient - unsatisfiable. A 250 W-at-25C part allows flange 75 C at 150 W -\textgreater{} Rth \textless{}= 0.333 C/W, which is a real heatsink. Must be re-derived from the actual datasheet curve at P1. \\
P0 & Match budget: residual series X \textless{}= 5.0 ohm (31.9 nH @ 25 MHz) for RL 26 dB & \textbar{}Gamma\textbar{} \textless{}= 0.0501 at 26 dB. Resistor +/-2\% DC tolerance costs only \textasciitilde{}0.1 dB, so launch/lead inductance is the entire design problem - drives the notch geometry, the tab bridge and the trimmer range. \\
P1 & BOM cost overrun accepted and reported, not engineered away & qty-5 BOM lands \textasciitilde{}4.3x over the \$40 cap. The cost is a materials floor (BeO/AlN RF ceramic, \textgreater{}=250W class), not a sourcing failure. Degrading to a 150W part or a wirewound would fail thermal feasibility or the 31.9 nH reactance budget - both hard specs. Report the real number (A8). \\
P1 & +/-5\% resistor tolerance rejected as a spec answer; re-check ordered & 'DC resistance 50 ohm +/-2\%' is a user-stated hard spec. The \textasciitilde{}0.1 dB RL argument is about return loss and is not permission to miss a separately-stated DC spec. If no \textless{}=2\% part is stocked, select-on-test is the fallback and gets documented as such. \\
P1 & 3 corner mounting holes (board\_only NPTH), giving exactly 6 total footprints & Budget is \textless{}=4 BOM lines / \textless{}=6 placements. Parts are R1+J1+C1 = 3 BOM lines, 3 CPL placements. Mounting holes are board\_only so they are not BOM lines (LEARNINGS 2026-07-29: a mounting hole is not a BOM line - fix at the symbol). Capping at 3 holes keeps TOTAL footprints at 6, so the constraint is met under the strict reading too. 3-point mount is also kinematically correct; 2 go at the SMA end where cable torque lands, 1 near the notch so the resistor tab solder joint is not a cantilever anchor. \\
P1 & DC 50 ohm +/-2\% met by SELECT-ON-TEST, not by datasheet tolerance & No in-stock \textgreater{}=250W flanged RF termination anywhere clears +/-2\%; T50R0's part-number scheme has no tolerance suffix at all, so tighter is pure special-order with no obtainable price. The board-level spec is therefore met at incoming inspection: 4-wire measure each part, accept 49.00-51.00 ohm. Documented in the README with the yield caveat - the datasheet only guarantees +/-5\%, so a build of 5 should order spares. \\
P1 & Ship BeO T50R0-250-12X; footprint also fits the AlN TA50R0-300-2X & The AlN twin is dimensionally identical (24.77 x 9.53 x 3.56 mm) but 0 stock at both distributors, and 'in stock' is a hard brief constraint. One footprint serves both, so the AlN part is a documented drop-in when stock returns. BeO toxicity gets a prominent handling warning in the README - do not machine, drill or break the part. \\
P1 & Keep BFC280811339 trimmer; Johanson 5602 recorded as alternate, not adopted & 9 pcs in stock covers a build of 5 and satisfies 'in stock'. The Johanson fixes stock depth but costs 7.1x (+\$272 on a build already 4.3x over budget) and is a different package, so it is not a drop-in. Bad value for a stock-depth risk; recorded as the substitute if stock lapses (footprint change required). \\
P2 & Shunt-C tuning cannot correct the resistive term - real limit is R\_eff = R + X\textasciicircum{}2/R & A shunt trimmer only cancels susceptance. Series residual X transforms the port resistance to R + X\textasciicircum{}2/R, which no amount of trimming undoes. At X=15.7 ohm (100 nH, the 'non-inductive' limit of ordinary bolt-down power resistors) R\_eff = 54.9 ohm -\textgreater{} 26.6 dB best case, and a +5\% high resistor drops that to 25.5 dB - FAILS the 26 dB spec. This is the quantitative reason an RF-rated low-L termination is mandatory and a cheap 0.1 uH thick-film part cannot substitute. \\
P2 & Required heatsink Rth = 0.633 C/W flange-to-ambient at 150 W / 25 C (0.500 C/W recommended) & Datasheet derating curve (T50R0-250-12X p.1) is 250 W flat to 100 C FLANGE, -5.0 W/C to zero at 150 C. 150 W is permitted at flange 120 C -\textgreater{} Rth = (120-25)/150 = 0.633 C/W. Staying at the 100 percent knee (100 C) needs 0.500 C/W and leaves 100 W of headroom. The curve is defined on flange temperature, so element-to-flange Rth (ABSENT from the datasheet) is not needed to close the requirement. \\
P2 & Interface resistance is a material share of the budget: flange is only 2.36 cm2 & A 0.5 C.cm2/W grease joint on a 24.77 x 9.525 mm flange is 0.212 C/W - a third of the whole 0.633 C/W budget. The heatsink itself must therefore be \textless{}= 0.42 C/W, not 0.63. Must be stated in the README or the number is misleading. \\
P2 & Resistor mounts OFF-BOARD beside the PCB; no notch. PCB rides on 1.0 mm shims & Tab underside is at 2.667 mm above the flange mounting plane; 1.6 mm PCB + 35 um copper puts top copper at 1.635 mm - a 1.03 mm gap. Three 1.0 mm shim washers lift top copper to 2.635 mm, leaving a 0.032 mm solder gap: an ideal flat lap joint with essentially zero added inductance. A notch clearing the full 24.77 x 9.525 mm flange would leave 2 mm rails on a 30 mm board (fragile), and the flange is 1.575 mm thick so it cannot tuck under a 1.0 mm-raised PCB anyway. Heatsink drilling template goes in the README. \\
P2 & Trimmer C1 goes in shunt at the SMA (port) end, not at the resistor end & The cancelling network is Y\_in = jwC + 1/(R + jwL): the shunt susceptance must be at the port side of the series inductance. Trimmer range 3-33 pF at 25 MHz cancels series X from 1.18 to 12.9 ohm = 7.5 to 82 nH. Below 7.5 nH the trimmer's own minimum over-corrects by \textless{}0.4 percent reflection (RL \textgreater{} 49 dB), so the low end costs nothing. Shunt C cannot cancel a CAPACITIVE residual - benign here since the residual is inductive by construction. \\
P2 & Board outline fixed at 24.0 x 16.0 mm (43 percent of the 30x30 HARD cap) & 16.0 axis is fully determined: 4.2 J1 pin inset (footprint copper half-extent 3.65 + 0.30 JLC edge rule) + 4.7 to the R1 pad (floor 4.43) + 6.6 lap pad + 0.5 edge. 24.0 axis is set by C1's dia 7.5 body on one side (12.65 min from the RF axis) and the M3 screw-head keepout on the other (8.35 min = 2.2 pad + 0.80 HV + 2.75 head radius + 2.6 hole-to-edge). True minimum is \textasciitilde{}21x15; the extra 3x1 mm is deliberate P6 packing margin - the outline is unrecoverable after P5 and JLC prices any 2L board flat below 100x100, so P6 freedom costs zero dollars and zero picohenries. \\
P2 & Launch length 4.7 mm; its FLOOR is 4.43 mm and is set by clearance, not by board size & The R1 lap pad is on the 122.5 Vpeak RF net and must clear J1's own ground pads by 0.80 mm: sqrt((d-2.55)\textasciicircum{}2 + 0.35\textasciicircum{}2) \textgreater{}= 1.90 gives d \textgreater{}= 4.42. Shrinking the board cannot shorten the launch. Rot 0 (RF exiting between two ground pads) is optimal; rotating J1 45 deg pushes the floor to 5.5 mm. Relaxing to row B4 (0.40 mm) would buy 0.9 mm = 0.4 nH = 0.06 ohm - not worth reopening a settled decision. \\
P2 & Residual series inductance estimated at 6.9 nH = 1.09 ohm at 25 MHz; RL 32.2 dB at the +5 percent corner & Term budget: 1.2 J1 internal + 0.40 J1 pad + 1.74 trace (3.55 mm at 1.1 mm wide, 0.489 nH/mm) + 1.50 lap pad (7.1 mm at 4.4 mm wide, 0.211 nH/mm) + 1.5 tab/gap/strap loop + 0.6 inside R1. Term F is BOUNDED not guessed: RL\textgreater{}=16 dB to 4 GHz implies X\textless{}=16.0 ohm there, so L\textless{}=0.64 nH. R\_eff = R + X\textasciicircum{}2/R gives 50.024 at R=50 (RL 73 dB), 51.023 at R=51 (RL 39.9), 52.522 at R=52.5 (RL 32.2). Spec is 26 dB. Failure needs X \textgreater{} 12.1 ohm = 77 nH, i.e. 11x margin. HEADLINE: the board meets 26 dB with the trimmer REMOVED (39.3 dB untuned at R=50); select-on-test buys 7.7 dB while the entire launch inductance costs 0.02 dB. \\
P2 & Launch is deliberately a LOW-Z wide line (4.4 mm, Z0 34 ohm), not 50 ohm; no impedance\_ohm declared & At 0.0017 lambda the structure is lumped, so the only figure of merit is L' = Z0{*}sqrt(eps\_eff)/c, which falls monotonically with width: 0.489 nH/mm at 1.1 mm, 0.303 at 2.83 mm (50 ohm), 0.211 at 4.4 mm. A 50 ohm launch would be 30 percent MORE inductive per mm. Extra shunt C lands mostly at the resistor end where it is nearly free. Also closes a trap: rules\_gen only solves impedance for DIFFERENTIAL pairs, so a single-ended impedance\_ohm would look declared and silently do nothing. \\
P2 & Two soldered copper ground straps bond the R1 flange top to F.Cu lands - the RF return must NOT go via the heatsink & The decided architecture covers the signal tab and is silent on the ground return, which is half the loop. Relying on the bolted flange-to-heatsink joint is quantifiably unsafe: 2.359 cm2 through 25 um of grease at er 5 is a 418 pF capacitor = 15.2 ohm at 25 MHz if the sink is anodised and has no asperity contact, plus a \textasciitilde{}60 mm2 detour loop worth 15-25 nH. That is a \textasciitilde{}10 dB board whose performance depends on a finish the USER picks. Two 0.3 x 4 x 3 mm straps give \textasciitilde{}0.3 nH instead. Sizing constraint: the element covers the flange's central 9.525 mm, so exposed flange top starts 4.76 mm off the tab axis - which is what forces the strap side of the board to 7.0 mm minimum. \\
P2 & C1 ROTOR to GND, stator to RF - binding on P3 pin mapping and P4 & requirements s8 F1: the operator turns the trimmer with a hand tool at 122.5 Vpeak while transmitting, and that is the intended use case, not a fault. The rotor is the terminal mechanically continuous with the adjustment screw, so grounding it puts the touched metal at ground potential. Stated fallback if the part does not distinguish the terminals: insulated tuning tool plus a README instruction to tune at reduced drive - declared explicitly, never picked silently. \\
P2 & Both copper layers poured GND (overrides the 2L B.Cu-only default); thermal key deliberately ABSENT & B.Cu is the microstrip return under the launch and is what makes L' 0.21-0.49 nH/mm instead of \textasciitilde{}0.8 nH/mm for a wire - it is why the budget closes at 6.9 nH. F.Cu is poured so the flange-bond lands and the C1 return exist without adding a footprint (the placement budget is exactly full at 6). No thermal key: A4 puts all 150 W through the flange off-board, so declaring it would make check\_thermal demand via arrays and copper for heat that never enters FR4, and place\_anneal's spreading term would fight the mandatory edge placement. The PCB dissipates \textless{} 0.25 W total. \\
P2 & ONE schematic sheet - P4 spawns a single schematic agent; nets are /RF (root local label) and GND (global) & 3 electrical parts, 6 footprints, 2 nets. Hierarchy would add sheet-prefix net names and a sheet-symbol layer to a circuit that fits on one line. A root-sheet local label spelled RF is named /RF, which is what every constraints.json entry spells; a hierarchical or global label instead makes check\_current, check\_creepage, check\_return\_path, rules\_gen's HV and width rules and planes\_gen's reference guarantee all silently pass on an empty set. Power symbols stay bare: GND. \\
P2 & Heatsink requirement verified against two real purchasable parts: 0.42 C/W IS reachable in free air, at 300x125x136 mm and 193 USD & Wakefield Thermal Solutions 392-300AB publishes 'Thermal Resistance @ Natural: 0.33 C/W' (DigiKey, 65 in stock, 193.02 USD qty 1, read 2026-08-08) -\textgreater{} Rth\_total 0.542 -\textgreater{} P = 625/(1+5{*}0.542) = 168 W, flange 106 C. Fischer Elektronik SK 88 100 SA (100x100x50 mm extrusion) publishes 1.05 C/W natural (DigiKey, 983 in stock, 11.85 USD) -\textgreater{} 85 W. Bare flange with no heatsink: 6.55 cm2 exposed at h=10-15 W/m2K -\textgreater{} 102-153 C/W -\textgreater{} \textasciitilde{}1 W, i.e. 0.7 percent of 150 W, so the README must say the assembly is unusable at rated power without the external sink. CAVEAT: both distributor fields say 'natural convection' but the delta-T/orientation fine print was unreadable (both vendors' PDFs refuse automated fetch), and RS lists the Fischer part at 0.85 K/W against DigiKey's 1.05 - unresolved. Ohmite AH50600V05000FE (1.30 C/W -\textgreater{} 73 W) is the one figure backed by a directly-read manufacturer PDF. \\
P2 & Mounting holes are PLATED, netted to GND via schematic MountingHole\_Pad symbols - they are RF circuit, not mechanical & R1's cold end returns through its flange to the heatsink, so the whole 1.732 Arms return is heatsink -\textgreater{} screws -\textgreater{} PCB GND plane -\textgreater{} J1 shell. board\_init --mounting-holes emits NPTH board\_only holes, which would leave the return path OPEN. Use 3x Mechanical MountingHole\_Pad symbols on GND with exclude\_from\_bom set at the SYMBOL (LEARNINGS 2026-07-29: a mounting hole is not a BOM line). Keeps 6 footprints / 3 BOM lines / 3 CPL placements. \\
P2 & Two of the three mounting screws must flank the R1 lap pad; shims must be METAL & Return-loop area sets the residual inductance. Bonding the ground plane to the heatsink right beside the resistor keeps the loop to roughly 6 mm x 2.6 mm (order 10 nH, \textasciitilde{}5 nH with two screws in parallel) instead of running the return the length of the board. The 1.0 mm shim washers are therefore part of the RF circuit and MUST be aluminium or brass - nylon or fibre washers open the return path. Goes in the README as an assembly requirement, not a suggestion. \\
P3 & C1 rotor to GND, stator to RF (adopting the librarian's recommendation) & Grounding the rotor puts the exposed adjustment screw at ground potential. That is both correct RF practice (the tuning tool and the operator's hand do not detune the reading) and a real safety improvement: the intended procedure is tuning while keyed, and the alternative leaves exposed metal at 122.5 Vpeak. \\
P2 & Adopt architect B3: soldered flange ground straps carry the RF return, not the mounting screws & Relying on the bolted flange-to-heatsink joint puts a 418 pF / 15.2 ohm series capacitor in the return at 25 MHz on an anodised sink - anodising is a dielectric, so the return would be open or wildly variable with the finish the user happens to pick. Two copper straps soldered from the silver-plated-copper flange to PCB GND lands make the return deterministic and short. Consequence: the mounting shims no longer need to be conductive, which frees them to be INSULATING and is now the main lever on trimmer temperature. \\
P2 & Outline 24.0 x 16.0 mm (43 percent of the 30x30 cap); launch length floor 4.42 mm & Architect: the 4.7 mm launch is set by the 0.80 mm clearance to J1's own ground pads, so shrinking the board cannot shorten it. Kept 3x1 mm packing margin because the outline is unrecoverable after P5 and JLC prices flat below 100x100. \\
P2 & Accept 2.18 mm guaranteed tab lap; pad 6.6 mm with a measure-and-trim build step & Datasheet gives tab protrusion as 3.18 mm MIN with no maximum, so the guaranteed lap after the mandatory 0.5 mm flange gap and 0.5 mm edge clearance is 2.18 mm, not the 3-5 mm assumed. 6.6 mm2 of lap solder carries 1.732 Arms comfortably; the tab is compliant 0.127 mm copper. Mitigation is a README incoming-inspection step, not a respin. \\
P2 & BOM/CPL accounting confirmed at 3 lines / 3 placements / 6 footprints & Assembly hardware (3 shims, 5 screws, 2 copper straps, grease) is user-supplied mounting hardware, not a PCB placement, and the brief scopes the heatsink and its attachment out of the design. Even the stricter reading that counts the straps as a BOM line gives 4, still inside the \textless{}=4 cap. \\
P3 & Trimmer = Johanson 5602 (DK 1956-1000-ND), 1-30 pF, 250 VDC, -65..+125 C, \$63.37/ea & BFC2808's whole family tops out at +85 C (PTFE sibling) and the PP part at +70 C, against a board that sits at 74-88 C. Chose 5602 over the 8052 because the brief requires a documented tuning PROCEDURE and the 5602's tool is confirmed (Johanson 8764, .130 in shank) while the 8052's is not - worth \$3.35 across the build. Also gives a wider documented adjustment range. Cost: trimmer goes \$44.42 -\textgreater{} \$316.85 per build of 5, now 72 percent of the BOM. \\
P3 & BFC280811339 retained as a documented reduced-duty alternative, not deleted & The +70 C problem only bites at full 150 W CW, and CW is the ORCHESTRATOR's assumption A2 - the user left the duty field blank. If the real application is pulsed at low duty the cheap part is fine and saves \$272 per build of 5. Keeping it as a priced, temperature-conditioned alternative puts that lever in the user's hands instead of silently spending their money. \\
P4 & Documented tuning range (Johanson 5602): exact null for residual L = 5.75 to 83.65 nH (X = 0.904 to 13.14 ohm at 25 MHz) & Trimmer 1-30 pF plus 1.3 pF pad parasitic = 2.3-31.3 pF total. Solved from the null condition wC = X/(R\textasciicircum{}2+X\textasciicircum{}2), not the L=CR\textasciicircum{}2 approximation, which is 7 percent optimistic at the top end. Low endpoint reproduces the sim-analyst's independent ngspice result (5.752 nH) to 4 digits. Separately, RL \textgreater{}= 26 dB is met from L = 0 up to \textasciitilde{}103 nH, because below the null floor the trimmer over-corrects by only a fraction of a pF. \\
P4 & Failure threshold corrected to X \textgreater{} 12.07 ohm / 76.9 nH at the +5 percent corner & sim-analyst solved R\_eff cap = 55.276 ohm for \textbar{}Gamma\textbar{} = 0.050119 and read Im\{Zin\} = 1.3e-6 ohm at the threshold, i.e. perfectly tuned with the residual reflection 100 percent real. The architect's 12.1 ohm is 0.2 percent optimistic - at exactly 12.1 ohm it reads 25.98 dB, marginally under spec. Quote 12.07. \\
P4 & Select-on-test, not launch inductance, is the entire error budget - confirmed by simulation & Corner sweep: a 3x change in inductance (5-\textgreater{}15 nH) moves RL by \textless{}=0.97 dB at fixed R, while the +5 percent resistance step costs 40.40 dB at fixed L. Worst 16-corner result is 31.90 dB, still 5.9 dB inside spec. This is the quantitative justification for spending the build effort on incoming resistance measurement rather than on layout inductance. \\
P3 & C1 net assignment defined by PHYSICAL terminal: case -\textgreater{} GND, insulated lead -\textgreater{} RF & No Johanson/Knowles document labels rotor vs stator for the 5602, and the librarian's inference chain had an acknowledged internal tension. Defining the netlist by physical terminal dissolves the question: a panel-mount trimmer's case is built to pass through a grounded panel, AND the case is the surface a hand or tool actually brushes. Grounding it makes 'exposed metal at ground potential' true regardless of which plate moves internally. No datasheet fact needed. \\
P3 & Tuning is performed at REDUCED drive or with a VNA, not at 150 W & The brief requires adjustment without desoldering or unbolting - it does not require adjustment at full power, and tuning for minimum reflection is a small-signal measurement that is more accurate at low drive anyway. This removes the live-at-122.5-Vpeak hazard from the normal procedure entirely, leaving the grounded case as defence in depth. Goes in the README procedure as a requirement, not a suggestion. \\
P3 & C1's 6.3 mm plated hole is AT JLC's standard limit (max\_hole\_diameter\_mm 6.3), not beyond it & jlc\_capabilities.yaml line 47. Compliant with the no-upcharge constraint, but with zero margin - flagged as a P9 DFM watch item rather than assumed safe. \\
P2 & R1 footprint and symbol reworked to 2 pins (RF tab + GND flange straps) & The first librarian pass built R1 as a single lap pad BEFORE the architect closed the return-path question. Block B3 needs two soldered flange-to-GND strap lands or there is physically nowhere to terminate the RF return, and a 1-pin symbol against a 2-pad footprint would fail P7 parity. \\
P5 & Outline grows from 24.0x16.0 to 26.0x20.0 mm before board\_init & As-built R1 lap pad is 7.0 mm vs blocks.md's 6.60 mm assumption, so the 16.00 mm launch axis was 0.40 mm over with zero margin. The outline binds permanently at board\_init and there is no shrink step, so the risk is asymmetric: too small costs a full P5-P7 redo, too large costs essentially nothing (JLC prices flat under 100x100, and the launch inductance penalty is \textasciitilde{}0.3 nH/mm against an 11x margin). 26x20 is 58 percent of the 30x30 HARD cap and leaves real room for C1's 7.6 mm pad, three M3 holes and 0.80 mm clearances. Extra area is SLACK - placement must still hold J1-to-R1 at the 4.42 mm clearance floor. \\
P9 & Heatsink needs a clearance hole under C1: \textgreater{}=7 mm dia, \textgreater{}=5 mm deep or through & C1's threaded case protrudes 4.75 mm below the PCB (6.35 mm stub minus 1.6 mm board) against 1.00 mm of shim clearance - a 3.75 mm interference. The README already owes a drilling template for the 2 flange bolts and 3 mounting holes, so one more clearance hole is a modest addition. Documented alternative: let the J1/C1 end of the board overhang the heatsink edge, which needs no machining AND lowers C1's temperature - both go in the README. \\
P3 & Canonical layer token is Dwgs.User, never User.Drawings & R1's first-pass footprint used the GUI display alias on 16 items. It round-trips fine through pcbnew and fp\_verify never sees it, but kicad-cli REFUSES to load the board (exit 3, no report) - it would have blocked every downstream gate with no diagnostic. Fixed and logged to LEARNINGS. \\
P5 & P5 complete: 26x20 2L JLC2313\_1.6, parity 0, 10 rules & board\_init parity cleared 3-\textgreater{}0 by adding Datasheet/Description properties to the LIBRARY footprints (incl. a local clone of J1's stock footprint) rather than by board surgery - the board was not yet placed or routed so regeneration was free. Root cause confirmed empirically: DRC parity compares only fields present on the schematic symbol INSTANCE, not inherited library defaults. \\
P9 & BLOCKING FOR P9: deliver the BOM from the schematic, not from bom\_cpl's pos export & Verified by reading bom\_cpl.py: build\_bom() iterates the parts list derived from the POS export, so --parts can only attach an LCSC to a ref already there - a footprint with exclude\_from\_pos\_files can never appear. R1 carries that attr (correctly - it is not a machine-placed part), so the delivered BOM would ship with C1 and J1 and NO 50 ohm load resistor. No gate catches this: dfm\_check's missing\_lcsc only inspects parts already in the pos report. Deliver BOM.csv authored from kicad-cli sch export bom (3 lines) with DigiKey PNs hand-carried, and use bom\_cpl only for the CPL. \\
P9 & CPL legitimately contains 2 rows (C1, J1); R1 is documented as not machine-placed & R1's body bolts to the user's heatsink off-board - its only board feature is the tab lap pad, so a CPL 'position' for it would mislead an assembler. Board is hand-built anyway (A5). README states BOM = 3 lines, CPL = 2 placements, both inside the \textless{}=4 line / \textless{}=6 placement cap. \\
P4 & H2 checkpoint auto-approved on the user's standing delegation & The user instructed: do not ask clarifying questions, make assumptions, record them, and proceed. Schematic reviewer returned 1 error (BOM completeness - a P9 delivery issue, not a circuit defect) and 3 warnings, all recorded with owners. The circuit itself reviewed CLEAN on every axis including the C1 safety property, verified in the committed .kicad\_pcb. Proceeding to P6; all findings surface in the final report. \\
P6 & Mounting-hole pad shrunk 6.4 -\textgreater{} 5.0 mm dia; outline stays 26x20 & Placement proved the intent geometrically impossible: legality uses pad-BBOX, making R1 a 14.2x7.5 mm solid rect and each M3 hole a 6.9 mm square, so only one hole fits in R1's Y-band. Flanking holes need offset \textgreater{}=10.55 mm but the board allows \textless{}=9.50 - short 1.05 mm per side. D\textless{}=5.35 closes it; D=5.0 gives a 0.9 mm annular ring on a 3.2 mm drill, well inside fab floors, and was verified legal with zero violations. Footprint swap only - growing the board to 28.1 mm would need a full P5 redo because the outline binds permanently. \\
P6 & OVERRIDING my own no-launch-trade instruction: move J1 to Yj=4.20, launch 4.50 -\textgreater{} 8.45 mm & New information the instruction predates: at the launch floor the SMA coupling nut sweeps to 0.45 mm from C1's body in plan, \textasciitilde{}2 mm above the board, on an axis height SCALED OFF A DRAWING and not published by the vendor. A nut that fouls the trimmer makes the board unusable and the tuning requirement unmeetable. Cost is +1.9 nH = 0.03 dB against 6.14 dB of margin and a 12.07 ohm / 76.9 nH failure threshold; it also removes 4 silk warnings. The instruction existed to stop casual spreading, not to override a mechanical conflict. \\
P6 & R1 footprint silk must be repaired: 0.12 mm text vs the 0.15 mm board floor, plus silk-over-copper & All 13 silk warnings belong to R1's own footprint graphics (BOLT CL / FLANGE/TAB CL labels), not to refdes, so place\_edit cannot fix them - it needs a footprint edit. drc\_routed requires err+warn = 0 at any severity, so these are blocking, not cosmetic. \\
P6 & AMENDS the P2 collinearity decision: J1 sits 2.70 mm west of R1's tab axis & Symmetric flanking holes pin R1's axis to x in {[}12.65,13.35{]}. With J1 collinear, C1 only fits rotated 90/270, which needs 11.1 mm of Y and leaves a 0.3 mm window where the M3 head keepout lands 0.64 mm from C1's RF pad - under the 0.80 mm HV floor. Rotating C1 to 180 needs 8.1 mm of Y but forces J1 to x \textless{}= 10.40. Cost of the offset is +0.39 mm of trace = +0.12 nH against a 9.5x margin. blocks.md s3.3 and constraints.json.\_note\_edges are superseded by this. \\
P6 & Flanking holes sit 3.87 mm north of the strap lands, not beside them & R1's own BOLT CL footprint silk (glyph top at y 16.695) vetoes anything lower. Benign under the adopted B3 decision because the SOLDERED STRAPS carry the RF return, not the screws - this geometry would only matter under the superseded screw-return assumption. \\
P8 & check\_silk treats an UNFILLED rect as solid ink - 5 false positives waived, tool bug logged & check\_silk.py builds Polygon(pts).buffer(w2) for any {*}\_rect and ignores the fill token. J1's F.SilkS rect is explicitly (fill no), 0.12 mm stroke. Proof: reported areas are exact whole-pad areas (3.80 mm2 = pi{*}1.1\textasciicircum{}2 corner pad; 5.26 vs the 2.3x2.3 centre pad), which only a filled interior produces; and kicad-cli DRC --severity-all reports zero silk findings on the same board while drc\_routed passes 0/0 ffailing on warnings too. Deforming a correct stock KiCad footprint to satisfy a checker that mis-parses fill would be the real defect. \\
P7 & Planes set to connect:solid rather than thermal relief & Thermal relief left J1's SE ground leg with 1 spoke against a minimum of 2 (starved\_thermal, a hard DRC error), and necking the RF could not fix it - even at the 0.9392 mm floor the surviving channel is 0.18 mm, under the 0.25 mm zone min\_thickness. Solid is also electrically correct: J1's legs carry the full 1.732 Arms return and 0.5 mm spokes are exactly the return-path slot that costs more than trace width buys. No thermal argument opposes it - all 150 W is off-board and the PCB dissipates under 0.25 W. \\
P7 & Router stopped flaring deliberately at 7.21 nH instead of minimising further & Flaring harder would push the residual BELOW C1's 5.75 nH exact-null floor, converting the trimmer's weak-but-real downward authority into no demonstrable null at all - i.e. optimising the headline inductance would have broken the brief's operator-adjustable requirement. 7.21 nH keeps a genuine two-sided null and is still 10.7x inside the 76.9 nH failure threshold. \\
P9 & Delivered BOM.csv authored from the schematic; bom\_cpl output would have shipped without the load resistor & Verified against the script: build\_bom() iterates the POS export, and R1 carries exclude\_from\_pos\_files (correctly - it is not machine-placed), so bom\_cpl produced 2 rows with R1 absent and C1 blank (DigiKey-only, and no script reads a DigiKey field). Delivered BOM built from kicad-cli sch export bom: 3 lines, distributor PNs hand-carried, qty-5 pricing, plus the user-supplied hardware list. Validated by parsing the CSV back and asserting all three refs present. \\
P9 & 6 silk strokes at 0.12 mm vs JLC's 0.15 mm minimum: documented, not fixed & Inherited from the stock KiCad SMA footprint. dfm gate PASSES (warning severity). Fixing needs surgery on a routed, gate-clean board plus a full drc\_routed/verify/dfm re-gate, to improve legibility of a connector outline whose orientation is already unambiguous from four symmetric ground legs and a centre pin. Regression risk outweighs a cosmetic gain. \\
\bottomrule
\end{longtable}
}
\subsection*{Sheet plan}
\subsection*{rf-term-150w - hierarchical sheet plan}
\subsubsection*{ONE SHEET. P4 SPAWNS A SINGLE SCHEMATIC AGENT.}
There is no hierarchy. The whole design is 3 electrical parts, 3 mechanical footprints and 2 nets. {*}{*}Everything lives on the root sheet; there are no child sheets and no hierarchical labels.{*}{*} Decomposing this would add sheet-prefix net names and a sheet-symbol layer of indirection to a circuit that fits on one line of text.

\begin{quote}\small\ttfamily\raggedright
\textbar{} Sheet \textbar{} Blocks (`blocks.md`) \textbar{} Refdes block \textbar{} `\#PWR` base \textbar{}\\
\textbar{}---\textbar{}---\textbar{}---\textbar{}---\textbar{}\\
\textbar{} {*}{*}root{*}{*} (`rf-term-150w.kicad\_sch`) \textbar{} B1 RF port, B2 adjustment, B3 termination element, B4 mechanical \textbar{} {*}{*}1-99{*}{*} (all prefixes) \textbar{} {*}{*}1{*}{*} \textbar{}
\end{quote}
Refdes uniqueness is trivial - one sheet, six refdes.

---

\subsubsection*{1. Net-naming contract - BINDING ON P4}
Only two nets exist. Getting their spelling wrong silently no-ops every P5-P8 consumer that reads \texttt{constraints.json}, so it is written here as a contract rather than left to chance.

\begin{quote}\small\ttfamily\raggedright
\textbar{} Canonical net \textbar{} How P4 must create it \textbar{} Carries \textbar{}\\
\textbar{}---\textbar{}---\textbar{}---\textbar{}\\
\textbar{} {*}{*}`/RF`{*}{*} \textbar{} a {*}{*}local label spelled `RF`{*}{*} placed on the root-sheet wire between J1.1, C1's stator pin and R1.1 \textbar{} 86.6 Vrms / 122.5 Vpeak / 1.732 Arms, DC-25 MHz \textbar{}\\
\textbar{} {*}{*}`GND`{*}{*} \textbar{} the {*}{*}global power symbol `GND`{*}{*} on J1.2, C1's rotor pin and R1.2 \textbar{} RF return + the flange bond \textbar{}
\end{quote}
Two rules behind that table:

1. \textbf{A root-sheet local label comes out as `/\textless{}LABEL\textgreater{}`, so `RF` becomes `/RF`.} Every \texttt{constraints.json} entry below spells it \texttt{/RF}. If P4 uses a hierarchical label, a global label or a different spelling, \texttt{check\_current}, \texttt{check\_creepage}, \texttt{check\_return\_path}, \texttt{rules\_gen}'s HV and width rules and \texttt{planes\_gen}'s reference guarantee all silently stop matching. There is no error - the checks just pass on an empty set. 2. \textbf{Power symbols are global and bare: `GND`, never `/GND`.}

\subsubsection*{2. Symbols, pins and the two pin-mapping calls P4 must get right}
\begin{quote}\small\ttfamily\raggedright
\textbar{} Ref \textbar{} Part class \textbar{} Symbol \textbar{} Footprint \textbar{} Pins -\textgreater{} nets \textbar{}\\
\textbar{}---\textbar{}---\textbar{}---\textbar{}---\textbar{}---\textbar{}\\
\textbar{} {*}{*}J1{*}{*} \textbar{} SMA female jack, right-angle THT (SMA-KWE class) \textbar{} `Connector:Conn\_Coaxial` \textbar{} `Connector\_Coaxial:SMA\_BAT\_Wireless\_BWSMA-KWE-Z001` \textbar{} 1 (centre) -\textgreater{} `/RF`, 2 (shield, 4 legs, one net) -\textgreater{} `GND` \textbar{}\\
\textbar{} {*}{*}R1{*}{*} \textbar{} 50 ohm 250 W flanged RF termination (T50R0-250-12X class) \textbar{} `Device:R` (2-pin passive) \textbar{} {*}{*}custom, P3 builds{*}{*} - see `blocks.md` s2.3 \textbar{} 1 (RF tab) -\textgreater{} `/RF`, 2 (flange, 2 lands) -\textgreater{} `GND` \textbar{}\\
\textbar{} {*}{*}C1{*}{*} \textbar{} 3-33 pF 250 V film trimmer, top-adjust (BFC2808 class) \textbar{} `Device:C\_Variable` \textbar{} {*}{*}custom, P3 builds{*}{*} - dia 7.5 mm THT radial, 2.5 mm pin grid \textbar{} {*}{*}stator -\textgreater{} `/RF`, ROTOR -\textgreater{} `GND`{*}{*} \textbar{}\\
\textbar{} {*}{*}H1-H3{*}{*} \textbar{} M3 mounting hole \textbar{} `Mechanical:MountingHole` \textbar{} `MountingHole:MountingHole\_3.2mm\_M3` \textbar{} none - {*}{*}`board\_only`{*}{*} \textbar{}
\end{quote}
\textbf{Call 1 - C1 rotor to GND, not to RF.} The rotor is the terminal mechanically continuous with the adjustment screw the operator touches with a tuning tool, at 122.5 Vpeak, while transmitting (the intended use case per A6). Grounding it puts the touched metal at ground potential. This is a requirements-s8-F1 safety consequence, not a preference. If the chosen part does not distinguish rotor from stator, say so and fall back to an insulated tuning tool plus a "tune at reduced drive" README instruction - \textbf{do not silently pick a pin.}

\textbf{Call 2 - H1-H3 must be `board\_only`.} A mounting hole that carries a symbol becomes a BOM line and a CPL line (LEARNINGS 2026-07-29 - fix it at the symbol, not downstream). The budget is \textless{}= 4 BOM lines and \textless{}= 6 placements; \texttt{board\_only} holes keep the BOM at 3 lines, the CPL at 3 placements and the total footprint count at exactly 6.

\subsubsection*{3. Interface nets and placement groups}
There are no inter-sheet interfaces. For P6, the two groupings that matter are:

\begin{quote}\small\ttfamily\raggedright
\textbar{} Group \textbar{} Members \textbar{} Rule \textbar{}\\
\textbar{}---\textbar{}---\textbar{}---\textbar{}\\
\textbar{} `port` \textbar{} anchor {*}{*}J1{*}{*}, member {*}{*}C1{*}{*} \textbar{} C1's shunt must tap the RF node at the port end. The tap point, not C1's body, is what matters - the stub between them is in series with C and shifts C\_eff by \textless{}0.3\% at 25 MHz (`blocks.md` s4.3). Keep the stub as short as C1's dia 7.5 mm body allows. \textbar{}\\
\textbar{} - \textbar{} {*}{*}C1{*}{*} vs {*}{*}R1{*}{*} \textbar{} `separation \textgreater{}= 10 mm`. Two independent reasons: the tuning-tool access cone must clear the flange (I2/A6), and C1 is a -40..+70 C part sitting near a 120 C flange (`blocks.md` s8 OPEN-1). \textbar{}
\end{quote}
Edge constraints (also in \texttt{constraints.json.placement.edges}):

\begin{quote}\small\ttfamily\raggedright
\textbar{} Ref \textbar{} Edge \textbar{} pos \textbar{} rot \textbar{} Why \textbar{}\\
\textbar{}---\textbar{}---\textbar{}---\textbar{}---\textbar{}---\textbar{}\\
\textbar{} J1 \textbar{} `top` \textbar{} 0.375 \textbar{} 0 \textbar{} at rot 0 this footprint's barrel and courtyard run to -12 mm in Y, so rot 0 on the top edge puts the 8.6 mm barrel off-board where a cable can mate. pos 0.375 = x 9.0 mm on a 24 mm board. \textbar{}\\
\textbar{} R1 \textbar{} `bottom` \textbar{} 0.375 \textbar{} 180 \textbar{} tab enters from off-board; flange, element and both mounting holes are outside the outline. \textbar{}
\end{quote}
\subsubsection*{4. ERC expectations}
\begin{itemize}
\item \textbf{Two nets, six footprints, zero power rails.} No \texttt{\#PWR} flag beyond \texttt{GND}.
\item \texttt{Device:R} and \texttt{Device:C\_Variable} are passive on both pins; \texttt{Conn\_Coaxial} is passive.
No driver anywhere - the board has no active part, so P4 should expect ERC to want a
\texttt{PWR\_FLAG} on \texttt{GND} (and only on \texttt{GND}).
\item No unconnected pins. No no-connects. No power-input pins.
\item \texttt{decoupling.json} is legitimately \textbf{empty/absent} - there is no IC to decouple
(\texttt{power\_tree.md} s1).
\end{itemize}
Full architecture narrative (not inlined; see the artifact index): \texttt{architecture/blocks.md}, \texttt{architecture/decisions.md}, \texttt{architecture/power\_tree.md}, \texttt{architecture/stackup.md}.
\section{Schematic}
Gate \texttt{erc} (phase P4): \textbf{pass} after 2 attempt(s); last 2026-08-09T02:12:27, failing 0 of 0.
The full schematic PDF follows.
\includepdf[pages=-]{reports/schematic.pdf}
\section{Layout}
Gate \texttt{place} (phase P6): \textbf{pass} after 1 attempt(s); last 2026-08-09T03:35:33, failing 0 of 0.

Gate \texttt{drc\_routed} (phase P7): \textbf{pass} after 1 attempt(s); last 2026-08-09T04:06:09, failing 0 of 0.
\begin{center}
\includegraphics[width=0.72\textwidth]{reports/renders/rf-term-150w_top.png}\\
{\small\texttt{reports/renders/rf-term-150w\_top.png}}
\end{center}
\begin{center}
\includegraphics[width=0.72\textwidth]{reports/renders/rf-term-150w_bottom.png}\\
{\small\texttt{reports/renders/rf-term-150w\_bottom.png}}
\end{center}
\section{Verification}
\emph{Pending --- produced at P8.}
\section{DFM and Fabrication}
\emph{Pending --- produced at P9.}
\section{Run Record (Appendix)}
\subsection*{Phase digests}
\paragraph*{P0}
\emph{(no digest recorded)}
\paragraph*{P1}
\emph{(no digest recorded)}
\paragraph*{P2}
\emph{(no digest recorded)}
\paragraph*{P3}
\emph{(no digest recorded)}
\paragraph*{P4}
\emph{(no digest recorded)}
\paragraph*{P5}
\emph{(no digest recorded)}
\paragraph*{P6}
\emph{(no digest recorded)}
\subsection*{Gate attempt history}
{\small
\begin{longtable}{llll}
\toprule
\textbf{Timestamp} & \textbf{Gate} & \textbf{Status} & \textbf{Failing/Total} \\
\midrule
\endhead
2026-08-09T00:33:55 & sim & pass & 0/0 \\
2026-08-09T01:47:45 & erc & pass & 0/0 \\
2026-08-09T02:12:27 & erc & pass & 0/0 \\
2026-08-09T03:35:33 & place & pass & 0/0 \\
2026-08-09T04:06:09 & drc\_routed & pass & 0/0 \\
2026-08-09T04:09:22 & verify & pass & 0/11 \\
2026-08-09T04:12:24 & dfm & pass & 0/1 \\
\bottomrule
\end{longtable}
}
\subsection*{Issues}
\emph{(no entries)}
\subsection*{Human checkpoints}
{\small
\begin{longtable}{lllp{2.6cm}p{6.4cm}}
\toprule
\textbf{Id} & \textbf{Phase} & \textbf{Status} & \textbf{When} & \textbf{Note} \\
\midrule
\endhead
1 & P2 & not reached &  &  \\
2 & P4 & approved & 2026-08-09T02:13:02 & auto-approved under the user's no-questions directive; reviewer: circuit clean, 1 P9 BOM-delivery error recorded, 3 warnings \\
3 & P6 & approved & 2026-08-09T03:36:21 & auto-approved under the no-questions directive; top render reviewed and matches intent; place gate 0 violations, DRC 0 warn / 0 parity \\
4 & P8 & not reached &  &  \\
5 & P10 & not reached &  &  \\
\bottomrule
\end{longtable}
}
\section{Artifact Index}
{\small
\begin{longtable}{p{9cm}rp{4cm}}
\toprule
\textbf{File} & \textbf{Size} & \textbf{Note} \\
\midrule
\endhead
\texttt{state.json} & 57,603 B &  \\
\texttt{kicad/rf-term-150w.kicad\_pcb} & 75,127 B &  \\
\texttt{kicad/rf-term-150w.kicad\_sch} & 38,840 B &  \\
\texttt{fab/rf-term-150w\_gerbers.zip} & 15,470 B &  \\
\texttt{fab/BOM.csv} & 2,610 B &  \\
\texttt{fab/CPL.csv} & 105 B &  \\
\texttt{brief/brief.md} & 2,235 B &  \\
\texttt{requirements.md} & 25,663 B &  \\
\texttt{architecture/blocks.md} & 32,809 B &  \\
\texttt{architecture/constraints.json} & 7,997 B &  \\
\texttt{architecture/decisions.md} & 9,234 B &  \\
\texttt{architecture/power\_tree.md} & 5,122 B &  \\
\texttt{architecture/sheets.md} & 5,180 B &  \\
\texttt{architecture/stackup.md} & 6,268 B &  \\
\texttt{reports/board\_init.json} & 9,642 B &  \\
\texttt{reports/bom\_cpl.json} & 1,225 B &  \\
\texttt{reports/CAP-ADJ-TH\_2P-BD7.5\_5602.overlay.svg} & 1,145 B &  \\
\texttt{reports/check\_return\_path.json} & 20,574 B &  \\
\texttt{reports/constraints\_lint.json} & 591 B &  \\
\texttt{reports/dfm\_check.json} & 1,201 B &  \\
\texttt{reports/fab\_export.json} & 2,330 B &  \\
\texttt{reports/gate-dfm.json} & 387 B &  \\
\texttt{reports/gate-drc\_routed.json} & 407 B &  \\
\texttt{reports/gate-erc.json} & 400 B &  \\
\texttt{reports/gate-place.json} & 337 B &  \\
\texttt{reports/gate-sim.json} & 364 B &  \\
\texttt{reports/gate-verify.json} & 14,805 B &  \\
\texttt{reports/netlist\_audit.json} & 962 B &  \\
\texttt{reports/p6-drc.json} & 8,658 B &  \\
\texttt{reports/p6-silk.json} & 1,098 B &  \\
\texttt{reports/place\_anneal.json} & 2,735 B &  \\
\texttt{reports/place\_metrics.json} & 2,733 B &  \\
\texttt{reports/place\_seed.json} & 1,720 B &  \\
\texttt{reports/plane\_repair.json} & 1,015 B &  \\
\texttt{reports/planes\_gen.json} & 828 B &  \\
\texttt{reports/R\_LapPad\_T50R0-250-12X.overlay.svg} & 1,493 B &  \\
\texttt{reports/review-schematic.json} & 5,868 B &  \\
\texttt{reports/review-schematic.md} & 15,789 B &  \\
\texttt{reports/route\_edit-rf.json} & 3,877 B &  \\
\texttt{reports/route\_edit-rmvias.json} & 757 B &  \\
\texttt{reports/rules\_gen.json} & 1,048 B &  \\
\texttt{reports/sch-bom.csv} & 429 B &  \\
\texttt{reports/schematic.pdf} & 57,200 B &  \\
\texttt{reports/silk\_place.json} & 1,596 B &  \\
\texttt{reports/sim.json} & 5,989 B &  \\
\texttt{reports/stitch\_vias.json} & 12,250 B &  \\
\texttt{reports/verify-waivers.json} & 3,580 B &  \\
\bottomrule
\end{longtable}
}
\end{document}
